\section{The dynamic-obstacle extension}
\label{sec:dynamic}

Static obstacles test avoidance; \emph{moving} obstacles test prediction. The
dynamic world adds three kinematic movers (two 0.7~m forklift-class boxes, one
1.0~m drum) driven by a shared motion contract, and the controller gains three
mechanisms to survive them.

\subsection{The shared motion contract}

Both the MATLAB simulation and the (future) ROS~2 stack evaluate obstacle
motion from one file, \code{dynamic\_obstacles.json}, interpreted by one
function (\code{eval\_mover\_pose.m}):
%
\[
\mathbf{m}(t) =
\begin{cases}
\mathbf{a} & t < t_0 \quad \text{(parked at the route start)}\\[2pt]
\mathbf{a} + \min\!\bigl(1,\, \tfrac{(t - t_0)\,v_m}{L}\bigr)(\mathbf{b} - \mathbf{a})
& \text{one-way}\\[4pt]
\text{shuttle } \mathbf{a} \leftrightarrow \mathbf{b},\ \text{leg time } 
T = \tfrac{L}{v_m} + t_{\mathrm{wait}} & \text{loop}
\end{cases}
\]
%
The velocity is piecewise constant, so \textbf{constant-velocity prediction is
exact between speed changes}, not an approximation. The scenario: mover1
crosses west$\to$east at $y = +0.5$ once (the formation yields); mover2 sweeps
the main aisle head-on forever; mover3 patrols the south aisle as background
traffic.

\subsection{Mechanism 1: predictive scan}

\code{pvfh\_predict.m} extends the VFH input: movers are raycast at their
contracted poses $t,\ t+\Delta t,\ \dots,\ t+T_{\mathrm{lead}}$
($T_{\mathrm{lead}} = 0.5$~s, sampled at $\Delta t$) and each beam keeps the
\emph{nearest blocking instance over the whole horizon}. The static sensing
path is untouched --- statics do not move, so their perception needs no
prediction. Effect: a mover closing head-on at $0.9$~m/s appears in the
histogram $\sim 0.5$~s ($\approx 0.45$~m of range) before an instantaneous
scan would place it inside the danger zone.

\begin{keybox}{A design dead end worth recording}
An earlier variant shifted the \emph{sensor position} by the self-motion
$\mathbf{v}_{\mathrm{self}} T_{\mathrm{lead}}$ and shortened all beams by the
self-displacement. That couples the current heading into the perception of
\emph{static} geometry: the controller sees different walls depending on which
way it points, and the feedback loop produces heading limit cycles. The fix
was architectural --- predict only the things that move. It is a good example
of why per-module numerics (scan construction) matter more than gains.
\end{keybox}

\subsection{Mechanism 2: speed slowdown}

The transport speed scales down smoothly as the \emph{predicted} worst-case
gap to any mover shrinks:
%
\begin{align*}
\text{speed\_scale} &=
\begin{cases}
\max\bigl(s_{\min},\ 0.8\, g_{\mathrm{pred}} / (m_{\mathrm{slow}} + r)\bigr)
& g_{\mathrm{pred}} < m_{\mathrm{slow}} + r\\
1 & \text{otherwise}
\end{cases}\\[6pt]
g_{\mathrm{pred}} &= \min_{T \le T_{\mathrm{look}}}
\Bigl[ \lVert \mathbf{m}(t{+}T) - \mathbf{c} - \mathbf{v}_{\mathrm{self}} T \rVert
- r_m - \lVert \mathbf{v}_{\mathrm{self}}\rVert T \Bigr]
\end{align*}
%
with $m_{\mathrm{slow}} = 0.25$~m, $s_{\min} = 0.4$, $T_{\mathrm{look}} =
1.0$~s. The floor $s_{\min}$ keeps the formation making progress (a full stop
near a moving obstacle is how deadlocks form); movers slower than 25\% of the
formation's speed are treated as parked and do not trigger the slowdown.

\subsection{Mechanism 3: the $t$ and $t{+}\Delta t$ commit test}

A step is committed only if the whole body (payload block included, via the
exact separating-axis test) is collision-free at $t$ \emph{and} at
$t + \Delta t$ with the movers at their contracted future poses. This is what
stops the block from stepping into a lane a forklift occupies one tick later.
The escape probe got the matching treatment: an escape heading is accepted
only if the body can hold it for $T_{\mathrm{hold}} = 1.5$~s while the movers
keep moving.

\subsection{Verified outcome}

(Single reference run recorded \emph{before} the October 2026 audit of
\S\ref{sec:audit}; re-run \code{Simulation8\_dynamic} to refresh it --- the
Monte Carlo tables of \S\ref{sec:experiments} are the authoritative
numbers.) On the dynamic world the full pipeline reaches the goal with
\textbf{zero collision steps}, minimum formation clearance $0.51$~m, minimum
gap to any mover $1.22$~m (the predicted worst case $0.33$~m, absorbed by the
slowdown and VFH), path efficiency $0.93$, total turning $12.0$~rad over
$42.5$~s of transport. The static-world regression on the same parameters is
unchanged, confirming the dynamic layer costs nothing when nothing moves.
